\documentclass[11pt,a4paper]{article}
\usepackage{curso}
\begin{document}
\portada{Lección}{10.7}{La arquitectura Z1T}

\noindent\textbf{Unidad 10.}\quad\textbf{Objetivos:} el modelo Z1T; la inferencia desagregada Z1 + FPGA; lectura crítica de sus tablas de energía y latencia.
\textbf{Prerrequisitos:} 9.5, 10.2--10.6.

\section{El modelo}
Z1T (Extropic, septiembre de 2026): $L=4$, $D=512$, $T=1024$. Atención convolucional con compuertas (GCA, familia AFT, 10.2), normalización DyT (10.3), MLP con unidades
tanh-lineales (10.6), proyecciones dispersas con $k=4$ entradas por salida (10.4), pesos de 4 bits y activaciones dy4p (\emph{dyadic}, 4 bits; 10.5). Cada pbit de salida
gasta sus 16 acoplamientos en 4 valores $\times$ 4 pbits dy4p. Según el artículo, igualar la pérdida de GPT-2 small exige del orden de diez veces más FLOPs de
entrenamiento.

\section{Inferencia desagregada}
Chip Z1: 8 núcleos, 269\,568 pbits, $>50$ MHz, $<1$ W; ejecuta las proyecciones muestreadas y DyT. FPGA: embedding, aritmética de GCA, residuos, codificación posicional,
cabeza de vocabulario. La FPGA podría sustituirse por GPU/XPU (no medido).

\section{Energía \evP}
$E_{\text{token}}=8{,}74\ (\text{Z1})+285{,}78\ (\text{FPGA})=294{,}52$ nJ: 97\,\% FPGA. H100 en decode con lote 1: 40{,}9 / 8{,}17 / 4{,}09 $\mu$J por token con MFU
10 / 50 / 100\,\%, es decir $\approx139\times$, $28\times$, $14\times$. Las capas de Z1 solas: $\approx4680\times$ / $935\times$ / $468\times$. El cuello de botella es la FPGA
(Amdahl, 9.5).

\section{Letra pequeña}
\begin{itemize}[nosep]
\item Cabeza de vocabulario excluida: con ella, $\approx136{,}4\ \mu$J por token ($\approx463\times$ el titular).
\item No se cuentan el movimiento de datos Z1--FPGA ni el número de chips; no se cuantizan las activaciones en la figura de calidad.
\item Latencia estimada $58{,}8\ \mu$s/token ($\approx17\,000$ tokens/s): 38{,}0 orquestación FPGA, 16{,}0 muestreo Z1 ($25\times32$ muestras a 50 MHz), 3{,}2 lectura,
1{,}3 escritura, 0{,}3 resto; H100 compilada: 102 $\mu$s.
\item Son proyecciones con lote 1, ancladas en experimentos con X0, no mediciones del sistema completo. Con lotes grandes la GPU sería mucho más eficiente.
\end{itemize}

\begin{aplicacion}
Receta de diseño; contabilidad por dispositivo con cabeza incluida; reducir el pegamento. Laboratorio L29: registro de mapeo y proyección de J/token.
\end{aplicacion}

\begin{error}
\begin{itemize}[nosep]
\item Citar $139\times$ sin el supuesto de MFU 10\,\%. \quad Ignorar la cabeza de vocabulario.
\item Confundir la razón de las capas Z1 con la del sistema. \quad Tratar proyecciones como mediciones.
\end{itemize}
\end{error}

\begin{thebibliography}{9}\small
\bibitem{z1t} Extropic, Z1T, \url{https://extropic.ai/writing/z1t} (septiembre de 2026).
\bibitem{aft} S. Zhai et al., An attention free transformer, arXiv:2105.14103 (2021).
\bibitem{dyt} J. Zhu et al., Transformers without normalization, \emph{CVPR} (2025).
\end{thebibliography}
\end{document}
